\scanpage{24}
\mspage{22}
% archive sheet number: 23
% hand: Grothendieck
% sheet_kind: cancelled_draft
% element: text | upper, lines 1–7 | Alors je dis que
Alors je dis que \deleted{\unsure{$E$}} \uline{si} $g\in G(R)$, \uline{$g$ est unipotent} ssi $g\in G^\circ(R)$. En effet, \deleted{\unsure{on est réduit|on est déjà}} à prouver \deleted{seulement ``seulement si''. \unsure{Sup}} seulement: $g$ unipotent $\Rightarrow g\in G^\circ(R)$. Or \deleted{\unsure{considérons son}} l'image de $g$ dans le groupe quotient $\Gamma=G/G^\circ$ est un \unsure{élém|élément} \uline{unipotent} de $\Gamma(R)$, et comme on est en car. $0$, c'est l'unité.

% element: text | cancelled middle, lines 7–9 | Pour toute
\deleted{\unsure{Pour toute}\illegible{2} Pour tout dire: on filtre $E$ par $E^0=0,\ldots E^i\ldots,\}$, \unsure{avec} \deleted{$E^i=$\unsure{image inverse de}}}
% element: display | cancelled middle, line 10 | E^(i+1) image inverse
\[
\deleted{E^{i+1}=\text{image inv. dans }E\text{ de }(E/E^i)^{G^\circ}.}
\tag*{(5.4.1)}
\]
% element: text | cancelled middle, lines 11–18 | Cette filtration
\deleted{Cette filtration est \deleted{\unsure{finie}}\inserted{exhaustive, i.e. $\exists i$ t.q. $E^i=E$,} ce qui exprime simplement que $G^\circ$ est unipotent. Alors $G/G^\circ$ opère sur les $E^{i+1}/E^i$, et si $g\in G$ est unipotent, il induit des op. unipotentes dans les $E^{i+1}/E^i$, donc des opérations triviales, i.e. il appartient au \uline{noyau} de $G/G^\circ$ opérant sur $\operatorname{Gr}(E)$. Mais comme on est en car. zéro, \deleted{\unsure{cela implique que}} l'op. de $G/G^\circ$ est \unsure{semi-simple}, donc $E\unsure{\sim|\simeq}\operatorname{Gr}(E)$ \unsure{par un}}

% element: text | lower, lines 19–22 | En particulier
En particulier, l'image inverse $U$ de $G^\circ(R)$ dans $\pi$ est le plus grand ss-groupe de $\pi$ opérant de façon unipotente, et à fortiori le plus grand ss-groupe ouvert ayant cette propriété.

% element: text | bottom, lines 23–27 | 5.5.a Notons
\unsure{5.5.a} Notons que \deleted{\unsure{l'on sait de}}\inserted{la structure} connue des \unsure{groupes alg.} \inserted{unipotents} commutatifs (sur un corps $R$ de car. $0$), comme groupes vectoriels, permet de \deleted{\unsure{faire}} trouver une filtration ``caractéristique'' de $G^\circ$ (par les \inserted{groupes de} commutateurs successifs, $G^{(0)}=G$,
